\chapter{Hydrodynamic Fluctuations And Long-Time Tails}
\label{ch:thermal-hydrodynamic-fluctuations}

\ConventionsLorentzian

Hydrodynamics is an effective theory of conserved densities.  At nonzero
temperature those densities fluctuate, and their nonlinear couplings produce
universal late-time contributions to correlation functions.  This chapter
derives the fluctuation framework and identifies the theorem boundary between
hydrodynamic effective theory and a microscopic QFT.

\section{Fluctuating Variables}

Let
\[
  \varphi_A(t,\mathbf x)
\]
denote hydrodynamic fields: energy density, charge densities, momentum
density, or equivalent variables such as temperature, chemical potentials,
and velocity.  Equilibrium fixes static susceptibilities
\[
  \chi_{AB}
  =
  \int d^{d-1}x\,
  \langle \delta q_A(0,\mathbf x)\delta q_B(0,\mathbf 0)\rangle_\beta .
\]
For a diffusive charge density \(n\), the linearized equation is
\[
  \partial_t n-D\nabla^2 n=\nabla_i\xi^i ,
\]
where the noise satisfies
\[
  \langle\xi^i(t,\mathbf x)\xi^j(t',\mathbf x')\rangle
  =
  2T\sigma\,\delta^{ij}\delta(t-t')\delta^{(d-1)}(\mathbf x-\mathbf x')
\]
and \(D=\sigma/\chi\).  The noise strength is fixed by the KMS condition,
equivalently by the fluctuation--dissipation relation.

\section{Schwinger--Keldysh Origin}

In the \(r,a\) basis, a quadratic diffusive effective action is
\[
  I[n_r,n_a]
  =
  \int dt\,d^{d-1}x\,
  \left[
    n_a(\partial_t n_r-D\nabla^2 n_r)
    +iT\sigma\,(\nabla_i n_a)(\nabla_i n_a)
  \right].
\]
Integrating over \(n_a\) gives the Langevin equation above.  The imaginary
part is positive because the microscopic density matrix defines positive
norms for Schwinger--Keldysh correlation functions.  Nonlinear hydrodynamics
adds interactions constrained by conservation laws, KMS symmetry, and
positivity.

\section{Mode-Coupling Integral}

Consider a conserved current containing a quadratic convective term
\[
  J^i=\lambda\,\varphi\,\varphi+\cdots .
\]
The one-loop hydrodynamic contribution to the retarded current correlator
contains an integral of the form
\[
  G_R^{JJ}(\omega)\supset
  \lambda^2
  \int\frac{d^{d-1}k}{(2\pi)^{d-1}}
  \frac{T^2\,\chi^{-2}}
       {-i\omega+2Dk^2}.
\]
At small \(\omega\), the nonanalytic part is
\[
  G_R^{JJ}(\omega)\sim
  \begin{cases}
    (-i\omega)^{(d-3)/2}, & d\ne3,\\
    \log(-i\omega), & d=3,
  \end{cases}
\]
up to coefficients fixed by the tensor structure and susceptibilities.  In
time domain this produces a long-time tail
\[
  C_{JJ}(t)\sim t^{-(d-1)/2}.
\]

\section{Cutoff Dependence And Transport}

The analytic part of the same loop integral depends on the hydrodynamic
cutoff \(\Lambda_{\mathrm{hyd}}\).  This dependence is absorbed into the
transport coefficient that appears in the constitutive relation at that
cutoff:
\[
  \sigma_{\mathrm{phys}}
  =
  \sigma(\Lambda_{\mathrm{hyd}})
  +
  \Delta\sigma(\Lambda_{\mathrm{hyd}}).
\]
The nonanalytic late-time contribution cannot be removed by a local
coefficient.  Therefore transport in finite-temperature QFT has two layers:
local coefficients in the hydrodynamic action and fluctuation corrections
computed within the effective theory.

\section{Stochastic Stress Tensor}

For a relativistic fluid, the stress tensor contains dissipative terms
\[
  T^{ij}
  =
  p\delta^{ij}
  +
  \rho v^iv^j
  -
  \eta
  \left(
    \partial^iv^j+\partial^jv^i
    -\frac{2}{d-1}\delta^{ij}\partial_kv^k
  \right)
  -
  \zeta\,\delta^{ij}\partial_kv^k
  +
  \Xi^{ij}.
\]
The noise satisfies
\[
\begin{aligned}
  \langle\Xi^{ij}(x)\Xi^{kl}(0)\rangle
  =
  2T\Bigl[
    \eta
    \left(
      \delta^{ik}\delta^{jl}+\delta^{il}\delta^{jk}
      -\frac{2}{d-1}\delta^{ij}\delta^{kl}
    \right)
    +\zeta\,\delta^{ij}\delta^{kl}
  \Bigr]\delta^{(d)}(x).
\end{aligned}
\]
This is the local equilibrium limit of stress-tensor fluctuation relations.
Its use requires a hydrodynamic cell scale larger than microscopic
correlation lengths and smaller than the external wavelength.

\section{Microscopic Theorem Boundary}

A derivation from QFT requires:
\[
\begin{gathered}
  \text{KMS state and conserved currents},\\
  \text{clustering and local equilibration},\\
  \text{existence of a hydrodynamic scaling limit},\\
  \text{positivity and KMS constraints on the Schwinger--Keldysh action},\\
  \text{control of the order of low-frequency and infinite-volume limits}.
\end{gathered}
\]
Hydrodynamic fluctuation theory computes consequences of these assumptions.
It does not by itself prove that a given microscopic QFT equilibrates.

\begin{openproblem}[Hydrodynamic limit]
\label{op:hydrodynamic-limit-qft}
Derive fluctuating hydrodynamics and long-time tails directly from an
interacting QFT with a KMS state, including the hydrodynamic scaling limit,
noise kernel, nonlinear couplings, and renormalization of transport
coefficients.
\end{openproblem}
